\section{Resultados}

\subsection{Validação cruzada agrupada}

A validação cruzada agrupada em cinco folds constitui o resultado primário do
estudo. No limiar 0,50 e sem TTA, a ConvNeXt-Tiny obteve AUC ROC média de
$0{,}9709\pm0{,}0059$ e AUC PR de $0{,}9607\pm0{,}0121$. A
Tabela~\ref{tab:cv-media} resume as métricas agregadas. Como a média pode ocultar
trocas operacionais entre falsos negativos e falsos positivos, a
Figura~\ref{fig:cv-folds} separa sensibilidade e especificidade por fold. A
função da figura é mostrar onde a validação cruzada varia: não no ranqueamento
global, mas no ponto de decisão binária usado em triagem.

\begin{table}[ht]
\centering
\caption{Resultado primário da validação cruzada em cinco folds agrupada por
imagem, com limiar 0,50 e sem TTA. Valores em média $\pm$ desvio-padrão.}
\label{tab:cv-media}
\begin{tabular}{lr}
\toprule
Métrica & Resultado \\
\midrule
AUC ROC & $0{,}9709\pm0{,}0059$ \\
AUC PR & $0{,}9607\pm0{,}0121$ \\
Acurácia balanceada & $0{,}9091\pm0{,}0115$ \\
Sensibilidade & $0{,}8994\pm0{,}0480$ \\
Especificidade & $0{,}9188\pm0{,}0400$ \\
\emph{F1} & $0{,}8926\pm0{,}0128$ \\
MCC & $0{,}8189\pm0{,}0208$ \\
Kappa & $0{,}8164\pm0{,}0223$ \\
\emph{Brier score} & $0{,}0699\pm0{,}0079$ \\
\bottomrule
\end{tabular}
\end{table}

\begin{figure}[ht]
\centering
\begin{tikzpicture}
\begin{axis}[
 width=.72\textwidth,height=4.6cm,
 kariellyaxis,
 ybar,bar width=9pt,
 symbolic x coords={1,2,3,4,5},xtick=data,
 xlabel={Fold},ylabel={M\'etrica dependente do limiar},
 ymin=.80,ymax=.99,
 legend style={at={(.5,-.20)},anchor=north,legend columns=3}
]
\addplot[fill=black!70,draw=black!70] table[x=fold,y=sensitivity,col sep=comma]{dados/cv_folds_plot.csv};
\addlegendentry{Sens.}
\addplot[fill=black!20,draw=black!45] table[x=fold,y=specificity,col sep=comma]{dados/cv_folds_plot.csv};
\addlegendentry{Espec.}
\addplot[black,thick,mark=*,sharp plot] table[x=fold,y=balanced_acc,col sep=comma]{dados/cv_folds_plot.csv};
\addlegendentry{Acc. bal.}
\end{axis}
\end{tikzpicture}
\caption{Métricas dependentes do limiar nos cinco folds agrupados. O fold 1
privilegia especificidade e perde sensibilidade; os folds 2 e 3 mostram o
comportamento oposto.}
\label{fig:cv-folds}
\end{figure}

A Figura~\ref{fig:cv-folds} justifica reportar a validação cruzada além da AUC:
no fold 1, a sensibilidade cai para 0,8253 enquanto a especificidade chega a
0,9681; nos folds 2 e 3, a recuperação de anormais aumenta, mas com menor
especificidade. Portanto, a variação relevante para uso em triagem está no
limiar operacional e na composição do fold, não apenas na capacidade de
ranqueamento do modelo.

\subsection{Análise detalhada em uma partição}

Na partição retida, a inferência da ConvNeXt-Tiny com TTA produziu AUC ROC de
0,9832 e AUC PR de 0,9827. A temperatura estimada na validação foi
$T=1{,}210$ e o \emph{Brier score} foi 0,0473. O ECE passou de 0,0247 antes da
transformação para 0,0130 depois dela, indicando melhora de calibração nessa
partição. Como o escalonamento por temperatura é monotônico, ele não altera as
AUCs; a pequena mudança de ranqueamento em relação à inferência sem TTA decorre
do TTA. As curvas ROC e precisão--revocação dessa partição estão na
Figura~\ref{fig:roc-pr}.

\begin{figure}[tbp]
\centering
\begin{tikzpicture}
\begin{groupplot}[
 group style={group size=2 by 1,horizontal sep=1.4cm},
 width=.40\textwidth,height=3.4cm,
 kariellyaxis,
]
\nextgroupplot[xlabel={1$-$Especificidade},ylabel={Sensibilidade},xmin=0,xmax=1,ymin=0,ymax=1,legend style={at={(.97,.06)},anchor=south east}]
\addplot[black,thick] coordinates {(0.0000,0.0000) (0.0000,0.1066) (0.0000,0.2102) (0.0000,0.3018) (0.0000,0.3944) (0.0000,0.4472) (0.0017,0.4980) (0.0026,0.5747) (0.0035,0.6315) (0.0052,0.6773) (0.0087,0.7221) (0.0105,0.7659) (0.0140,0.8038) (0.0166,0.8416) (0.0245,0.8735) (0.0359,0.9014) (0.0481,0.9293) (0.0674,0.9482) (0.0945,0.9582) (0.1207,0.9691) (0.1522,0.9761) (0.1855,0.9821) (0.2178,0.9871) (0.2511,0.9900) (0.2878,0.9920) (0.3263,0.9940) (0.3902,0.9950) (0.4576,0.9950) (0.5372,0.9950) (0.5914,0.9970) (0.6938,0.9970) (0.7822,0.9990) (0.8959,0.9990) (0.9563,1.0000) (1.0000,1.0000)};
\addlegendentry{AUC 0,9832}
\addplot[black!45,dashed] coordinates {(0,0) (1,1)};
\nextgroupplot[xlabel={Revoca\c{c}\~ao},ylabel={Precis\~ao},xmin=0,xmax=1,ymin=.44,ymax=1,legend style={at={(.03,.06)},anchor=south west}]
\addplot[black,thick] coordinates {(0.0010,1.0000) (0.1096,1.0000) (0.2171,1.0000) (0.3257,1.0000) (0.4343,1.0000) (0.4602,0.9978) (0.5129,0.9961) (0.5667,0.9965) (0.6195,0.9936) (0.6713,0.9912) (0.7211,0.9864) (0.7729,0.9835) (0.8217,0.9786) (0.8675,0.9710) (0.9054,0.9548) (0.9213,0.9448) (0.9392,0.9374) (0.9502,0.9235) (0.9582,0.9075) (0.9691,0.8726) (0.9741,0.8564) (0.9821,0.8244) (0.9871,0.7922) (0.9920,0.7477) (0.9940,0.7058) (0.9950,0.6682) (0.9950,0.6335) (0.9960,0.6031) (0.9970,0.5753) (0.9980,0.5422) (0.9990,0.5197) (0.9990,0.4921) (1.0000,0.4797) (1.0000,0.4676)};
\addlegendentry{AUC 0,9827}
\addplot[black!45,dashed] coordinates {(0,0.4676) (1,0.4676)};
\end{groupplot}
\end{tikzpicture}
\caption{Curvas ROC e precisão--revocação da ConvNeXt-Tiny na partição retida
(com TTA e temperatura). A linha tracejada indica o classificador aleatório
(ROC) e a prevalência de células anormais (precisão--revocação).}
\label{fig:roc-pr}
\end{figure}

\subsection{Efeito do limiar operacional}

Os três regimes evidenciam a troca entre falsos negativos e falsos positivos. O
limiar 0,50 apresentou o maior MCC e a maior acurácia balanceada. O limiar de
Youden aumentou a sensibilidade com perda moderada de especificidade. O regime
de alta sensibilidade atingiu 0,9701 de sensibilidade, mas reduziu a
especificidade para 0,8723. Em um fluxo de triagem, esse ponto poderia ser
aceitável caso a prioridade fosse reduzir omissões, transferindo mais casos para
revisão humana.

Os intervalos de confiança por bootstrap agrupado por imagem reforçam que a
diferença entre 0,50 e Youden é pequena na partição retida. Para o limiar 0,50,
a sensibilidade foi 0,9402 (IC95\% 0,8801--0,9737), a especificidade foi 0,9440
(IC95\% 0,9087--0,9697) e o MCC foi 0,8840 (IC95\% 0,8203--0,9261). Para Youden,
a sensibilidade foi 0,9522 (IC95\% 0,8939--0,9829), a especificidade foi 0,9283
(IC95\% 0,8925--0,9561) e o MCC foi 0,8791 (IC95\% 0,8150--0,9218). Para o limiar
de alta sensibilidade, a sensibilidade foi 0,9701 (IC95\% 0,9342--0,9902), a
especificidade foi 0,8723 (IC95\% 0,8239--0,9139) e o MCC foi 0,8414 (IC95\%
0,7766--0,8892).

\subsection{Comparação com ResNet50}

A comparação interna com ResNet50 foi realizada na mesma partição retida. A
ConvNeXt-Tiny apresentou maior capacidade de ranqueamento, com AUC ROC 0,9816 e
AUC PR 0,9805, contra 0,9710 e 0,9750 da ResNet50. A ResNet50, por outro lado,
obteve maior especificidade, acurácia balanceada e MCC no limiar avaliado. Esse
resultado sugere que a ConvNeXt-Tiny ordena melhor as amostras, enquanto a
decisão binária depende fortemente do limiar escolhido.

\begin{table}[ht]
\centering
\caption{Comparação entre ConvNeXt-Tiny e ResNet50 na partição retida.}
\label{tab:baseline}
\kariellytablefont
\kariellytablecols
\begin{tabular}{lrrrrrrr}
\toprule
Modelo & AUC ROC & AUC PR & Acc. bal. & Sens. & Espec. & F1 & MCC \\
\midrule
ConvNeXt-Tiny & \textbf{0,9816} & \textbf{0,9805} & 0,9371 & \textbf{0,9373} &
0,9370 & 0,9331 & 0,8738 \\
ResNet50 & 0,9710 & 0,9750 & \textbf{0,9401} & 0,9064 & \textbf{0,9738} &
\textbf{0,9362} & \textbf{0,8852} \\
\bottomrule
\end{tabular}
\end{table}

\subsection{Representatividade e erros por categoria Bethesda}

A composição da partição retida favorece categorias mais facilmente separáveis:
HSIL corresponde a 52,7\% das células anormais no teste, ante 35,8\% no
conjunto completo, enquanto LSIL corresponde a 16,4\%, ante 28,6\% no conjunto
completo. Essa super-representação de HSIL e sub-representação de LSIL ajudam a
explicar por que os resultados da partição única são superiores às médias da
validação cruzada.

No limiar de Youden, a Tabela~\ref{tab:bethesda} mostra 92,8\% de acerto para
NILM, 81,7\% para ASC-US, 94,5\% para LSIL, 90,2\% para ASC-H, 99,4\% para HSIL
e 100\% para SCC. O resultado de ASC-US, baseado em 115 células dessa partição,
sugere uma dificuldade relevante, mas não sustenta sozinho uma conclusão
populacional. SCC também possui suporte reduzido, com 42 células.

\begin{table}[ht]
\centering
\caption{Acerto binário por categoria Bethesda na partição retida, usando o
limiar de Youden (0,287).}
\label{tab:bethesda}
\small
\setlength{\tabcolsep}{5pt}
\begin{tabular}{lrrlr}
\toprule
Categoria & $n$ & Acerto (\%) & IC95\% Wilson & $P$(anormal) média \\
\midrule
NILM & 1143 & 92,8 & 91,2--94,2\% & 0,091 \\
ASC-US & 115 & 81,7 & 73,7--87,7\% & 0,741 \\
LSIL & 165 & 94,5 & 90,0--97,1\% & 0,898 \\
ASC-H & 153 & 90,2 & 84,5--94,0\% & 0,860 \\
HSIL & 529 & 99,4 & 98,3--99,8\% & 0,956 \\
SCC & 42 & 100,0 & 91,6--100,0\% & 0,970 \\
\bottomrule
\end{tabular}
\end{table}
